%%%%%%%%%%%%%%%%%%%%%%%%%PORTADA Y DATOS%%%%%%%%%%%%%%%%%%%%%%%%%%%%%
\begin{center}
\textbf{ACTIVIDAD AUTÓNOMA} \\
\vspace{0.5cm}
\textbf{Cifrado Criptográfico: Implementación de Mejores Prácticas en el Sistema de Gestión Nutricional IMC}
\end{center}

\begin{center}
\begin{tabular}{|l|l|}
\hline
   \textbf{Estudiante}  &  Kevin Patricio Sarango Olaya \hspace{3.5cm} \\
\hline
   \textbf{Paralelo} &  8VO \hspace{3.5cm} \\
\hline
   \textbf{Professor} &  Ing. Cristian Narváez G. Mg.Sc.\\
\hline
   \textbf{Unidad} & Unidad 2: Cifrado y Protección de Datos \\
\hline
   \textbf{Fecha} & 30/04/2026 \hspace{1cm} \textbf{Plagio:} Sí:\quad No\quad \hspace{2.4cm} \textbf{Calificación:} ..... \\
\hline
   \multicolumn{2}{|l|}{\textbf{Observación.:} } \\
\hline
\end{tabular}
\end{center}

%%%%%%%%%%%%%%%%%%%%%%%%%%%%%%%%%%TEXTO%%%%%%%%%%%%%%%%%%%%%
\section*{Detalles de la Actividad}
\section*{Pregunta de investigación: ¿Es suficiente con cifrar los datos, o la solidez de una
aplicación depende de la correcta selección de algoritmos, la gestión segura de llaves y el uso
de sales aleatorias?}

\section*{Resultados de Aprendizaje}
\begin{enumerate}
    \item \textbf{Investigar} y seleccionar algoritmos de cifrado simétrico y asimétrico adecuados
          para el contexto del proyecto integrador.
    \item \textbf{Implementar} mecanismos de protección de datos en tránsito (TLS~1.3) y datos en
          reposo (base de datos y sesiones).
    \item \textbf{Validar} la robustez criptográfica mediante la métrica de Resistencia Criptográfica
          y el análisis de resistencia a ataques conocidos.
\end{enumerate}
